\section{Conclusion}
We studied learning-to-defer for non-stationary time series under partial feedback and
time-varying expert availability. We proposed \textbf{L2D-SLDS}, a factorized switching
linear-Gaussian state-space model over signed expert residuals with context-dependent regime
transitions, a shared global factor, and per-expert idiosyncratic states. On a synthetic
regime-dependent correlation benchmark, on Melbourne Daily Temperatures, and on Jena Climate,
L2D-SLDS consistently improves over contextual bandit baselines, non-stationary bandit baselines
(D-LinUCB, CUSUM-LinUCB, GLR-LinUCB), and over ablations of the shared factor, regime count,
transition model, and EM estimation. A registry-churn experiment further shows that the maintained
online state scales with the active registry rather than the cumulative expert catalog.

\textbf{Limitations.}
The paper does not provide regret-style or minimax guarantees; the theoretical contribution is
structural (cross-expert information transfer, pruning invariance, coupling at birth) rather than
regret-theoretic. Deriving meaningful regret bounds in a setting that simultaneously combines
non-stationary latent regimes, partial feedback, expert heterogeneity, and time-varying availability
remains an open direction. Additionally, routing quality depends on the quality of the SLDS
parameters, whether set manually or estimated via EM; the model is a hierarchy where practitioners
select the complexity that fits their setting, and simpler operating points (\(M=1\), uniform
transitions, no EM) are available as special cases.


\section{Impact Statement}
This paper aims to advance the field of machine learning. While the proposed methods may have broader societal implications, we do not identify any specific societal consequences that require explicit discussion at this stage.
